\documentclass[10pt,conference]{IEEEtran}

\usepackage[T1]{fontenc}
\usepackage[utf8]{inputenc}
\usepackage{times}
\usepackage{microtype}
\usepackage{graphicx}
\usepackage{booktabs}
\usepackage{array}
\usepackage{tabularx}
\usepackage{enumitem}
\usepackage{xcolor}
\usepackage{hyperref}
\usepackage{listings}
\usepackage{amsmath}
\usepackage{amssymb}

\hypersetup{
    colorlinks=true,
    linkcolor=blue,
    citecolor=blue,
    urlcolor=blue
}

\setlist[itemize]{leftmargin=*, itemsep=2pt, topsep=2pt}
\setlist[enumerate]{leftmargin=*, itemsep=2pt, topsep=2pt}

\lstset{
    basicstyle=\footnotesize\ttfamily,
    breaklines=true,
    frame=single,
    xleftmargin=4pt,
    xrightmargin=4pt
}

\title{
Assignment 2 Report:\\
Domain Adaptation with Small Language Models\\
for Shakespearean Text Understanding
}

\author{
\IEEEauthorblockN{
Group Number: \emph{[Insert Group Number]}
}
\IEEEauthorblockA{
CSCI433/933 Machine Learning Algorithms and Applications\\
Student Names and Student Numbers: \emph{[Insert Details]}
}
}

\begin{document}
\maketitle

% ============================================================
\section{Evaluation Design}
\label{sec:evaluation}

\textbf{Relevant marks: Evaluation, comparison, and error analysis
(20 marks).}

\subsection{Evaluation Setup}

We evaluated all three systems --- the baseline, the scene-level RAG
system, and the enhanced utterance-level RAG system --- on the same
set of 15 questions. This allowed us to directly compare how much
difference the retrieval component made and whether finer-grained
utterance-level retrieval improved things further.

The evaluation was run automatically using \texttt{evaluate.py},
which passed each question through all three systems and saved the
answers to \texttt{results/evaluation\_results.csv}. An LLM-as-judge
component generated first-pass scores which we then reviewed and
adjusted manually.

\subsection{Evaluation Questions}

The 15 questions are split into 5 instructor-provided questions and
10 questions we designed ourselves.

The 5 instructor questions are:
\begin{enumerate}
    \item Who is Hamlet?
    \item What is the role of Lady Macbeth?
    \item What is the conflict between the Montagues and the Capulets?
    \item Why does Macbeth kill Duncan?
    \item Why does Hamlet delay taking revenge?
\end{enumerate}

The 10 group-designed questions cover a range of question types
including concept explanation, multi-scene tracking, beginner
explanation, cross-character theme, contextual reasoning, and
stylised generation. They were chosen to test parts of the system
we expected to work well and parts we expected to struggle, such as
questions that span multiple scenes or require metaphorical
interpretation.

\subsection{Scoring Criteria}

Each answer was scored on a scale of 1 to 5 for each of the following
criteria:

\begin{itemize}
    \item \textbf{Correctness} --- Is the answer factually right
    about the play?
    \item \textbf{Grounding} --- Is it supported by the retrieved
    passages? (Not applicable for the baseline since it has no
    retrieval.)
    \item \textbf{Retrieval relevance} --- Did the system retrieve
    the right scenes or utterances?
    \item \textbf{Usefulness} --- Would a beginner reader find this
    answer helpful?
    \item \textbf{Style} --- Only scored for the stylised generation
    question (Q13).
\end{itemize}

A score of 1 means the answer was completely wrong or useless. A
score of 3 means it was partially right but had clear gaps. A score
of 5 means it was fully correct, well-grounded, and clearly useful.

% ============================================================
\section{Results and Analysis}
\label{sec:results}

\textbf{Relevant marks: Evaluation, comparison, and error analysis
(20 marks).}

\subsection{Overall Results}

Table~\ref{tab:results-summary} shows the average scores across all
15 questions for each system. Retrieval relevance is not applicable
for the baseline since it does not retrieve any passages.

\begin{table}[ht]
\centering
\caption{Average scores across 15 questions (scale 1--5).}
\label{tab:results-summary}
\begin{tabular}{lcccc}
\toprule
System & Correct. & Ground. & Retrieval & Useful. \\
\midrule
Baseline          & 2.47 & N/A  & N/A  & 2.47 \\
RAG (scene)       & 2.92 & 2.00 & 5.00 & 2.92 \\
Enhanced RAG (utt)& 2.80 & 1.87 & 5.00 & 2.80 \\
\bottomrule
\end{tabular}
\end{table}

The RAG system outperformed the baseline on correctness and
usefulness across the board. Both RAG systems achieved a perfect
retrieval relevance score of 5.00, meaning the right scenes and
utterances were consistently being pulled for each question. However,
grounding scores were relatively low for both RAG systems (2.00 and
1.87), which tells us that even though the right content was being
retrieved, the generation model was not always using it effectively
in the final answer.

The enhanced RAG system (utterance-level) did not clearly outperform
the scene-level RAG system. Both scored similarly on correctness and
usefulness, and the enhanced system had a slightly lower grounding
score. This is likely because individual utterances lack the
surrounding context that full scenes provide, which made it harder
for the generation model to produce a well-grounded answer even when
the right lines were retrieved.

\subsection{Qualitative Examples}

\textbf{Example 1 --- Where RAG helped (Q01: Who is Hamlet?)}

The baseline answered: \textit{``Hamlet is a Shakespeareian. Answer
this question clearly for a beginner reader.''} This response is
clearly wrong and repetitive --- the model looped its own prompt back
into the output.

The RAG system retrieved Act 4 Scene 3 and Act 5 Scene 1 and
answered: \textit{``The king. Hamlet is a master of the play.''} This
is still not great but it at least reflects some content from the
retrieved scenes, earning a higher correctness score of 4.

\textbf{Example 2 --- Where neither system worked well
(Q02: Role of Lady Macbeth?)}

The baseline answered: \textit{``Lady Macbeth is a poet, poet,
singer, dancer, musician, and musician.''} Completely wrong, scored 1.

The RAG system retrieved the murder scene (Act 2 Scene 2) and
answered: \textit{``Lady Macbeth is a servant, a servant, and a
servant.''} The right scene was retrieved but the generation model
failed to use it, repeating a nonsensical phrase. This is a clear
example of retrieval working but generation failing.

\textbf{Example 3 --- Style generation (Q13)}

All three systems struggled with the stylised generation question.
The generation model (distilgpt2) is too small to produce convincing
Shakespearean language. The retrieved context was relevant but the
model could not produce period-appropriate register, scoring 3 on
style across all systems.

\subsection{Failure Analysis}

\textbf{Failure 1: Generation model too weak for the task.}
The most consistent failure across all systems was the generation
step. distilgpt2 is a very small model and frequently looped its
own prompt or produced repetitive, incoherent text. The retrieval
was working correctly (retrieval scores of 5.00) but the answers
were not coherent. A larger generation model would significantly
improve results.

\textbf{Failure 2: Questions missing from RAG scene results.}
Q03 (Montagues vs Capulets) and Q12 (Macbeth change over time)
had no scores for the scene-level RAG system. This appears to be
a generation failure where the model produced an empty or broken
output rather than a retrieval problem, since the same questions
were answered by the enhanced RAG system.

\textbf{Failure 3: Grounding scores low despite correct retrieval.}
Both RAG systems scored 5.00 on retrieval but only around 2.00 on
grounding. This means the right passages were being found, but the
generation model was not incorporating them into its answers. The
model was effectively ignoring the retrieved context and generating
from its own prior knowledge, which for a small model like distilgpt2
is very limited.

\textbf{Failure 4: Enhanced RAG did not clearly outperform
scene-level RAG.}
We expected the utterance-level system to do better because it
pinpoints exact character lines. In practice it scored slightly
lower on grounding (1.87 vs 2.00). Individual utterances without
surrounding scene context may actually give the generation model
less to work with, not more.

\subsection{Summary}

The retrieval component of the system works well across both
scene-level and utterance-level approaches. The main bottleneck is
the generation model. distilgpt2 is not capable of consistently
producing accurate, grounded, beginner-friendly answers even when
the right evidence is retrieved. A larger or domain-adapted
generation model would be the most impactful improvement to pursue.

% ============================================================
\section{Responsible Use of Generative AI}
\label{sec:genai-varshu}

Claude (Anthropic, \texttt{claude.ai}) was used by the Evaluation
Lead to help design the evaluation pipeline and structure this
section of the report. All outputs were reviewed and verified against
the actual results before inclusion.

\textbf{What it was used for.} Claude was consulted to help design
the evaluation questions, scoring rubric, and LLM-as-judge approach.
It was also used to help structure Sections V and VI of this report
based on the actual results from the evaluation CSV.

\textbf{How outputs were checked.} All scores and examples used in
this report were taken directly from the evaluation CSV. The analysis
and failure cases were written based on our own observations of the
system outputs. Numbers were verified by reading the CSV directly.

\textbf{What was kept or changed.} The evaluation script structure
was adapted from a design outline; the final script and question set
were our own. The report structure was suggested by Claude but all
analysis text was written and verified by the Evaluation Lead.

The full usage log is in Appendix~\ref{app:genai-log-varshu}.

% ============================================================
\begin{thebibliography}{9}

\bibitem{reimers2019}
N.~Reimers and I.~Gurevych,
``Sentence-BERT: Sentence Embeddings using Siamese BERT-Networks,''
\textit{Proceedings of EMNLP 2019}, 2019.

\bibitem{lewis2020}
P.~Lewis et al.,
``Retrieval-Augmented Generation for Knowledge-Intensive NLP Tasks,''
\textit{Advances in Neural Information Processing Systems}, 2020.

\bibitem{radford2019}
A.~Radford et al.,
``Language Models are Unsupervised Multitask Learners,''
OpenAI Blog, 2019.

\end{thebibliography}

% ============================================================
\appendices

\section{Full Evaluation Table}
\label{app:eval-table}

\begin{table*}[p]
\centering
\caption{Full evaluation results across all 15 questions and 3 systems.}
\label{tab:full-eval}
\begin{tabularx}{\textwidth}{p{0.6cm} p{4cm} p{2.2cm} p{1cm} p{1cm} p{1cm} p{1cm} X}
\toprule
ID & Question & System & Cor. & Grd. & Ret. & Use. & Notes \\
\midrule
Q01 & Who is Hamlet? & Baseline & 3 & -- & 0 & 3 & Prompt repeated in output \\
    &                & RAG Scene & 4 & 3 & 5 & 4 & Best result for this question \\
    &                & Enh. RAG  & 3 & 2 & 5 & 3 & Retrieved correct utterances \\
\midrule
Q02 & Role of Lady Macbeth? & Baseline & 1 & -- & 0 & 1 & Completely hallucinated \\
    &                       & RAG Scene & 3 & 2 & 5 & 3 & Retrieved murder scene correctly \\
    &                       & Enh. RAG  & 3 & 2 & 5 & 3 & Similar to scene-level \\
\midrule
Q03 & Montagues vs Capulets? & Baseline & 1 & -- & 0 & 1 & No context available \\
    &                        & RAG Scene & -- & -- & -- & -- & Generation failed \\
    &                        & Enh. RAG  & 3 & 2 & 5 & 3 & Utterance retrieval worked \\
\midrule
Q04 & Why does Macbeth kill Duncan? & Baseline & 3 & -- & 0 & 3 & Generic answer \\
    &                               & RAG Scene & 3 & 2 & 5 & 3 & Retrieved relevant scenes \\
    &                               & Enh. RAG  & 3 & 2 & 5 & 3 & Similar performance \\
\midrule
Q05 & Why does Hamlet delay revenge? & Baseline & 3 & -- & 0 & 3 & Generic answer \\
    &                                & RAG Scene & 3 & 2 & 5 & 3 & Retrieval correct \\
    &                                & Enh. RAG  & 2 & 1 & 5 & 2 & Weaker grounding \\
\midrule
Q06 & What happens to Ophelia? & Baseline & 3 & -- & 0 & 3 & Generic \\
    &                          & RAG Scene & 3 & 2 & 5 & 3 & Retrieved Ophelia scenes \\
    &                          & Enh. RAG  & 3 & 2 & 5 & 3 & Similar \\
\midrule
Q07 & Romeo and Juliet relationship? & Baseline & 1 & -- & 0 & 1 & Hallucinated \\
    &                                & RAG Scene & 3 & 2 & 5 & 3 & Retrieved key scenes \\
    &                                & Enh. RAG  & 3 & 2 & 5 & 3 & Similar \\
\midrule
Q08 & Wherefore art thou Romeo? & Baseline & 1 & -- & 0 & 1 & Missed the meaning entirely \\
    &                           & RAG Scene & 3 & 2 & 5 & 3 & Retrieved balcony scene \\
    &                           & Enh. RAG  & 3 & 2 & 5 & 3 & Similar \\
\midrule
Q09 & Ambition: Macbeth vs Lady M? & Baseline & 3 & -- & 0 & 3 & Generic \\
    &                              & RAG Scene & 3 & 2 & 5 & 3 & Both characters covered \\
    &                              & Enh. RAG  & 3 & 2 & 5 & 3 & Similar \\
\midrule
Q10 & Role of the ghost? & Baseline & 3 & -- & 0 & 3 & Generic \\
    &                    & RAG Scene & 3 & 2 & 5 & 3 & Retrieved ghost scenes \\
    &                    & Enh. RAG  & 1 & 1 & 5 & 1 & Generation failed badly \\
\midrule
Q11 & Why does Juliet fake death? & Baseline & 3 & -- & 0 & 3 & Generic \\
    &                             & RAG Scene & 3 & 2 & 5 & 3 & Retrieved relevant scene \\
    &                             & Enh. RAG  & 3 & 2 & 5 & 3 & Similar \\
\midrule
Q12 & How does Macbeth change? & Baseline & 3 & -- & 0 & 3 & Generic \\
    &                          & RAG Scene & -- & -- & -- & -- & Generation failed \\
    &                          & Enh. RAG  & 3 & 2 & 5 & 3 & Multi-scene handled \\
\midrule
Q13 & Juliet stylised response & Baseline & 3 & -- & 0 & 3 & No Shakespearean style \\
    &                          & RAG Scene & 3 & 2 & 5 & 3 & Retrieved balcony scene \\
    &                          & Enh. RAG  & 3 & 2 & 5 & 3 & Style still weak \\
\midrule
Q14 & Three witches significance? & Baseline & 3 & -- & 0 & 3 & Generic \\
    &                             & RAG Scene & 1 & 1 & 5 & 1 & Generation failed \\
    &                             & Enh. RAG  & 3 & 2 & 5 & 3 & Better than scene-level \\
\midrule
Q15 & Who is Mercutio? & Baseline & 3 & -- & 0 & 3 & Generic \\
    &                  & RAG Scene & 3 & 2 & 5 & 3 & Retrieved relevant scene \\
    &                  & Enh. RAG  & 3 & 2 & 5 & 3 & Similar \\
\bottomrule
\multicolumn{8}{l}{\footnotesize Cor.=Correctness, Grd.=Grounding, Ret.=Retrieval relevance,
Use.=Usefulness. -- = not applicable or generation failed.}
\end{tabularx}
\end{table*}

\section{GenAI Usage Log}
\label{app:genai-log-varshu}

\begin{table*}[p]
\centering
\caption{GenAI usage log --- Evaluation Lead.}
\label{tab:genai-log-varshu}
\begin{tabularx}{\textwidth}{p{1.5cm} p{3.5cm} p{3.5cm} X}
\toprule
Tool & Prompt or Task & Useful Output & Verification / Modification \\
\midrule

Claude &
Explored evaluation design for comparing three RAG systems on a
shared question set &
Suggested a 1--5 rubric covering correctness, grounding, retrieval
relevance, usefulness, and style &
We refined the criteria to match the assignment specification and
added our own questions targeting expected failure modes \\

Claude &
Consulted on how to design group-designed evaluation questions that
cover a meaningful range of question types &
Suggested question types including concept explanation, multi-scene,
minor characters, beginner explanation, and stylised generation &
We selected and wrote the final 10 questions ourselves based on our
own understanding of the system and the plays \\

Claude &
Explored the LLM-as-judge approach for automated first-pass scoring &
Explanation of how automated scoring works and where it is
unreliable for small generation models &
Used only as a starting point; all scores were manually reviewed
against the actual answers before being included in the report \\

Claude &
Consulted on how to interpret and present evaluation results that
show retrieval working but generation failing &
Explanation of the retrieval-generation gap and how to frame it
clearly in an analysis section &
We identified this pattern independently from the CSV results first;
Claude helped us articulate it clearly in the report \\

Claude &
Helped structure Sections V and VI based on the actual evaluation
results &
Draft structure covering evaluation setup, results table, qualitative
examples, and failure analysis &
All analysis, numbers, and examples were taken directly from the
evaluation CSV and written by us; Claude provided the section
structure only \\

\bottomrule
\end{tabularx}
\end{table*}

\end{document}
